\chapter{Lattice Coefficients and Moment Transform}
\label{app:coefficients}

\section{Velocity sets and weights}
\label{sec:app-weights}

The D3Q27 velocity set used for the momentum population groups into four
families by squared length. Table~\ref{tab:d3q27-weights} gives the weights of
the equilibrium~\eqref{eq:equilibrium} for each family.

\begin{table}[h]
\centering
\caption[D3Q27 and D3Q7 lattice weights]{Lattice weights for the momentum
population on D3Q27 and the scalar population on D3Q7. Velocities are given in
units of $\Delta x / \Delta t$.}
\label{tab:d3q27-weights}
\begin{tabular}{llrrl}
\toprule
Set & Family & $|\mathbf{c}_i|^{2}$ & Count & Weight \\
\midrule
\multirow{4}{*}{D3Q27}
 & rest      & 0 & 1  & $8/27$ \\
 & face      & 1 & 6  & $2/27$ \\
 & edge      & 2 & 12 & $1/54$ \\
 & corner    & 3 & 8  & $1/216$ \\
\midrule
\multirow{2}{*}{D3Q7}
 & rest      & 0 & 1  & $1/4$ \\
 & face      & 1 & 6  & $1/8$ \\
\bottomrule
\end{tabular}
\end{table}

\section{Moment ordering and relaxation rates}
\label{sec:app-moments}

The transform $\mathbf{M}$ of Equation~\eqref{eq:mrt-collision} is built by
Gram Schmidt orthogonalisation of the polynomial basis
$\{1, c_x, c_y, c_z, c^{2}, 3c_x^{2}-c^{2}, c_y^{2}-c_z^{2}, c_x c_y, c_y c_z,
c_x c_z, \dots\}$ against the weighted inner product
$\langle a, b \rangle = \sum_i w_i a_i b_i$. The production relaxation rates
are given in Table~\ref{tab:relaxation-rates}. Rates for moments above fourth
order are set to unity, which is the maximally dissipative choice and is
appropriate because those moments carry no hydrodynamic information.

\begin{table}[h]
\centering
\caption[Relaxation rates by moment order]{Relaxation rates used in the
production configuration. The shear rate $s_\nu$ is set from the target
viscosity through Equation~\eqref{eq:viscosity} and is the only rate that varies
between refinement levels.}
\label{tab:relaxation-rates}
\begin{tabular}{llll}
\toprule
Moment group & Order & Symbol & Rate \\
\midrule
Density, momentum   & 0, 1 & $s_\rho, s_j$ & 1.0 (conserved) \\
Deviatoric stress   & 2    & $s_\nu$       & from~\eqref{eq:viscosity} \\
Bulk stress         & 2    & $s_e$         & 1.19 \\
Heat flux           & 3    & $s_q$         & 1.40 \\
Fourth order        & 4    & $s_4$         & 1.72 \\
Ghost               & 5, 6 & $s_g$         & 1.00 \\
\bottomrule
\end{tabular}
\end{table}

The value $s_e = 1.19$ is inherited from the linear stability analysis
of~\citet{manov2017mrt} and was verified here against the Taylor Green decay of
Section~\ref{sec:val-order}, where a value below $1.1$ produced a visible
oscillation in the bulk mode and a value above $1.4$ over damped the acoustic
transient during spin up.
